Um die neue Wellenlänge $\lambda^{\prime}$ eines Photons nach dem Compton-Streuungsvorgang zu bestimmen, verwenden wir die Compton-Gleichung. Diese beschreibt die Änderung der Wellenlänge eines Photons, das an einem ruhenden Teilchen der Masse $m$ gestreut wird, in Abhängigkeit vom Streuwinkel $\theta$:

\begin{align*}
\lambda^{\prime} - \lambda &= \frac{h}{m c} (1 - \cos \theta).
\end{align*}

Hierbei ist $h$ das Plancksche Wirkungsquantum und $c$ die Lichtgeschwindigkeit. Die Gleichung zeigt, dass die Wellenlängenänderung proportional zur Compton-Wellenlänge $\frac{h}{m c}$ des Teilchens ist und vom Streuwinkel $\theta$ abhängt. 

Um die neue Wellenlänge $\lambda^{\prime}$ zu bestimmen, lösen wir die Gleichung nach $\lambda^{\prime}$ auf:

\begin{align*}
\lambda^{\prime} &= \lambda + \frac{h}{m c} (1 - \cos \theta).
\end{align*}

Diese Gleichung beschreibt, wie die Wellenlänge des Photons nach der Streuung zunimmt, was als Compton-Verschiebung bekannt ist.

%CHECK: Die Lösung war korrekt und vollständig; keine Schwächen oder Fix-Suggestions im Review.